\section{一般位置六点における corner の Hilbert 関数の一様性}
\label{sec:hilbert-constant}
\begin{theorem}[corner の完全な Hilbert 関数]\label{thm:hilbert-constant}
標数零の代数閉体上，十八の係数 $w_{1bc},w_{2bc}$ が非零で，
定理\ref{thm:general-position-iff}の十一の非零条件を満たすとする．
$C=B/BeB$ に道長 grading を入れると，六点配置には依存せず
\[
 \boxed{H_C(t)=8+21t+33t^2+38t^3+33t^4+21t^5+8t^6,\qquad
       \dim_k C=162.}
\]
特に定理\ref{thm:uniform-c7}の $\dim C_6\ge8$ は等号であり，
$C_0,\ldots,C_6$ の次元跳びは一般位置 locus に存在しない．
この結論は $B$ 全体が AS 正則であることを仮定しない．
\end{theorem}

\begin{proof}
まず定理\ref{thm:general-position-iff}により，
$C$ を一般位置の六点を吹き上げた滑らかな可換三次曲面 $S$ の
rolled-up helix 代数の同じ idempotent 商として計算できる．
既存の記号を保ち，$H=-K_S$，
\[
\begin{gathered}
 D_1=-2h+\sum_{r=1}^6e_r,\qquad D_2=-h+e_1+e_2+e_3,\\
 A_0=\mathcal O_S,\quad A_a=\mathcal O_S(D_a)\ (a=1,2),\quad
 B_i=\mathcal O_S(e_i)\ (1\le i\le3),\\
 C_j=\mathcal O_S(G_j),\qquad G_j=h-e_j\ (4\le j\le6),\qquad
 q_j=H-G_j=2h-\sum_{r\ne j}e_r.
\end{gathered}
\]
ここで $q_j$ は五点を通る conic の真の変換である $(-1)$ 曲線の類を表す．
$A,B,C$ の三 block にレベル $0,1,2$ を与えると，
道長 $n$ の始点 $s$，終点 $t$ の全 Hom は，
AOU\cite[(4.1), (4.9)]{AOU}の同定により
\begin{equation}\label{eq:hilbert-block}
 B_n(s,t)=H^0\!\left(S,\mathcal O_S\left(
 \left\lfloor\frac{\ell(s)+n}{3}\right\rfloor H+E_t-E_s\right)\right)
 \quad(\ell(t)\equiv\ell(s)+n\pmod3),
\end{equation}
ただし $E_{A_a}=D_a$，$E_{B_i}=e_i$，$E_{C_j}=G_j$ である．
AOU の helix の acyclicity と曲面の Riemann--Roch
$\chi(\mathcal O(D))=1+D(D+H)/2$ を使うと，
$e=A_0$ を除く八頂点の全 Hom の各次数の合計は次の表の中央左列になる．
\[
\begin{array}{c|rrrrrrrr}
\text{道長 }n&0&1&2&3&4&5&6&7\\\hline
\dim B_n\text{（八頂点間）}&8&21&42&74&105&147&206&252\\
\dim (BeB)_n\text{（同上）}&0&0&9&36&72&126&198&252\\
\dim C_n&8&21&33&38&33&21&8&0
\end{array}
\]
以下では中央右列，すなわち構造頂点を\emph{途中で経由する}
積の像の次元を各次数について証明する．
多重経由があるときは，すべての経由位置に対応する像の和を取る．

必要な切断数を列挙する（$a=1,2$，$i=1,2,3$，$j=4,5,6$）：
\begin{equation}\label{eq:hilbert-basic-h0}
\begin{gathered}
 h^0(H)=4,\quad h^0(2H)=10,\quad h^0(H\pm D_a)=3,\\
 h^0(H-e_i)=2,\quad h^0(H+e_i)=5,\quad h^0(2H-e_i)=7,\\
 h^0(G_j)=2,\quad h^0(q_j)=1,\quad
 h^0(H+G_j)=7,\quad h^0(H+q_j)=5.
\end{gathered}
\end{equation}
$H-e_i$ と $G_j$ は base-point-free pencil である．
$H\pm D_a$ は定理\ref{thm:uniform-c7}で使った平面への blow-down 線束である．

\paragraph{道長二と三．}
道長二では $C_j\to B_i$ の九組のみが $e$ を経由し，
その像は $H^0(q_j)H^0(e_i)$ の一次元である．合計九次元．
道長三では $A_a\to A_b$ に経由射はなく，
$B_i\to B_k$ と $C_j\to C_k$ の九組ずつで，
それぞれ $H^0(H-e_i)H^0(e_k)$ および
$H^0(q_j)H^0(G_k)$ が二次元の像を与える．合計三十六次元．

\paragraph{道長四．}
全二十一組の全 Hom はそれぞれ五次元である．
$A_a\to B_i$（六組）では $H^0(H-D_a)H^0(e_i)$ の三次元が経由する．
$B_i\to C_j$（九組）では
$H^0(H-e_i)H^0(G_j)$ が経由する．後者の pencil trick の核は
$H^0(q_j-e_i)$ であり，これは零である．実際 $q_j\cdot e_i=1$ と
$q_j^2=-1$ から，そのような effective divisor は $q_j$ を成分として含むが，
差は $-e_i$ となって effective でない．従って積の像は四次元．
$C_j\to A_a$（六組）では $H^0(q_j)H^0(H+D_a)$ の三次元が経由する．
この三種類を加えると $6\cdot3+9\cdot4+6\cdot3=72$ である．

\paragraph{道長五．}
全二十一組の全 Hom はそれぞれ七次元である．
$A_a\to C_j$（六組）では
$H^0(H-D_a)H^0(G_j)$ の像を取る．$G_j$ の pencil trick の核は
$H^0(H-D_a-G_j)$ であり，零である．
$a=1$ なら $H-D_1-G_j=2q_j-e_j$ で，
$q_j$ との交点数は $-2$ だから $q_j$ が二重の固定成分となり，
残りが $-e_j$ となって矛盾する．$a=2$ なら
$H-D_2-G_j=3h-2(e_1+e_2+e_3)-\sum_{r=4,r\ne j}^6e_r$．
これは互いに素な三本の $(-1)$ 曲線
$h-e_1-e_2$，$h-e_1-e_3$，$h-e_2-e_3$ とそれぞれ交点数 $-1$ であり，
三本を引くと $-\sum_{r=4,r\ne j}^6 e_r$ となり effective でない．
よって積の像は $3\cdot2=6$ 次元．
$B_i\to A_a$（六組）では
$H^0(H-e_i)H^0(H+D_a)$ の像を取る．$H-e_i$ の pencil trick の核は
$H^0(D_a+e_i)=0$（$h$ との交点が負）なので六次元．
$C_j\to B_i$（九組）は $e$ を二通りの位置で経由し，
$H^0(H+q_j+e_i)$ の中で
\[
 q_j\,H^0(H+e_i)+e_i\,H^0(H+q_j)
\]
を得る．ここで曲線の記号はその定義切断を掛ける作用も表す．
二つの部分空間は共に五次元であり，$q_j$ と $e_i$ に共通成分がないので
交叉は $q_je_i\,H^0(H)$ の四次元である．よって和は六次元．
従って経由射の合計は $(6+6+9)\cdot6=126$ 次元である．

\paragraph{道長六：$A_a\to A_b$（四組）．}
この場合の経由射は一種類で，次の積の像である：
\begin{equation}\label{eq:two-blowdowns-injection}
 H^0(H-D_a)\otimes H^0(H+D_b)
 \longrightarrow H^0(2H+D_b-D_a).
\end{equation}
この積写像が全ての $a,b=1,2$ に対して単射であることを示す．
$L=H-D_a$，$M=H+D_b$ とする．いずれも $S\to\mathbb P^2$ の
blow-down を定め，切断空間は三次元である．
$L$ が縮約する六本の $(-1)$ 曲線のうち，$M\cdot E=2$ となる曲線は
$a=b$ で六本，$a\ne b$ で三本（後者は $q_4,q_5,q_6$）ある．
$L$ により縮約されるそれら三本の像は $\mathbb P^2$ で非共線である．
もし共線なら，その直線の真の変換は $H$ と交点零の有効曲線となり，
$H$ の ampleness に矛盾する．
一方，$M\cdot E=2$ なら，$M$ による blow-down は $E$ を
滑らかな平面 conic に同型に写す．したがって
\[
 H^0(M)\xrightarrow{\sim}H^0(E,M|_E)\simeq H^0(\mathbb P^1,\mathcal O(2))
\]
である．
仮に \eqref{eq:two-blowdowns-injection} の核に $\sum_{u=1}^3s_u\otimes t_u$
（$s_u$ は $H^0(L)$ の基底）があれば，三本の曲線への制限から
$\sum_{u=1}^3s_u(p_E)t_u=0$ が $H^0(M)$ で成り立つ．
三つの $p_E$ が非共線なので，三つの評価汎関数は独立であり，
全ての $t_u$ が零である．これで積写像の単射性が従う．
従って全ての四組で経由射の像は九次元である．

\paragraph{道長六：$B_i\to B_k$（九組）．}
$H^0(H-e_i)H^0(H+e_k)$ の像は，
$H-e_i$ の pencil trick と $h^0(e_i+e_k)=1$ より九次元である．
$i\ne k$ では全 Hom が九次元なのでこれで終わる．
$i=k$ のときは二番目の経由位置から $e_i\,H^0(2H-e_i)$ も得られるが，
$H^0(H-e_i)H^0(H)\twoheadrightarrow H^0(2H-e_i)$ が成立する．
実際，pencil trick の核は $H^0(e_i)$ の一次元で，
$2\cdot4-1=7=h^0(2H-e_i)$ だからである．
従って二番目の積は最初の像に含まれ，九次元のままである．

\paragraph{道長六：$C_j\to C_k$（九組）．}
$H^0(H+q_j)H^0(G_k)$ の像は，$G_k$ の pencil trick と
$h^0(q_j+q_k)=1$ より九次元である．
$j\ne k$ では全 Hom が九次元なのでこれで終わる．
$j=k$ の場合，もう一つの経由位置から $q_j\,H^0(H+G_j)$ も得るが，
$H^0(H)H^0(G_j)\twoheadrightarrow H^0(H+G_j)$ が成立する．
実際，pencil trick の核は $H^0(q_j)$ の一次元で，
$4\cdot2-1=7=h^0(H+G_j)$ である．
従って後者の積も先の九次元の像に含まれる．

以上より道長六の経由射の次元は
$(4+9+9)\cdot9=198$ である．
表中の全 Hom の次元 $206$ との差は八であり，$\dim C_6=8$ が従う．
道長七以降の消滅は定理\ref{thm:uniform-c7}による．
表を合計して Hilbert 関数と $\dim C=162$ を得る．
\end{proof}

\begin{corollary}[162次元 algebra の有限平坦族]
\label{cor:flat-corner-162}
標数零の代数閉体 $k$ 上の四交比トーラスにおける十一条件の補開集合
$U=\{\Delta_{11}\ne0\}$ 上，規格化した universal quiver potential の
$e_{A_0}$-corner 商を $\mathcal C_U$ とする．このとき
$\mathcal C_U$ は階数 $162$ の有限局所自由な
$\mathcal O_U$-代数で，その道長斉次部分は順に
\[
 (8,21,33,38,33,21,8)
\]
の階数を持ち，道長七以上は零である．任意の $U$ 上の基底変更と可換する．
\end{corollary}
\begin{proof}
各道長 $n$ の成分は有限個の道を基底とする自明ベクトル束を，
全ての $u(\partial W)v$ からなる有限個の関係によって割った
有限表示 $\mathcal O_U$ 加群である．
定理\ref{thm:hilbert-constant}より $0\le n\le6$ では
全幾何学的ファイバーの次元が一定であるから，
対応する関係行列は一定階数を持ち，商は局所自由で基底変更と可換する．
$n=7$ では全ファイバーが零なので，Nakayamaの補題で
$\mathcal C_{U,7}=0$．頂点と矢による生成性から高次数も零になる．
有限個の斉次部分を足し合わせて主張を得る．
\end{proof}

\begin{remark}[結果の対象範囲]\label{rem:hilbert-status}
証明は四交比の十一条件による一般位置六点の再構成と，
可換三次曲面の幾何学的 rolled-up helix の Hom 同定に基づく．
$w_{0bc}$ は $C$ に影響しないので任意でよい．
しかし $27$ 係数の任意の potential が AS-regular/noetherian であること，
Rains heart や VdB heart と同一であること，導来圏の全域での同値は
ここからは結論しない．
\end{remark}
